\documentclass[10pt]{article}
\usepackage[preprint]{tmlr}
\usepackage{amsmath,amssymb}
\usepackage{booktabs}
\usepackage{enumitem}
\usepackage{graphicx}
\usepackage{hyperref}
\usepackage{url}
\hypersetup{hidelinks}

\def\month{09}
\def\year{2026}
\def\openreview{\url{https://openreview.net/forum?id=XXXX}}

\title{What a Benign-Trained Anomaly Detector Can and Cannot See:\\
A Real-Data Study of the Masquerade--Traitor Boundary on TWOS}

\author{\name Suramya Angdembay \email suramya.angdembay@usm.edu \\
\addr University of Southern Mississippi
\AND
\name Haiyan Tian \email haiyan.tian@usm.edu \\
\addr University of Southern Mississippi}

\begin{document}
\maketitle

\begin{abstract}
Benign-trained anomaly detectors are evaluated by how well their scores rank
attacks, but a score says nothing about \emph{which} attacks a detector can
see at all. We study this directly on TWOS, a dataset of $24$ real users with
genuine self-reported personality inventories and two distinct kinds of staged
insider activity---\emph{masquerades}, in which an attacker operates a
victim's live session, and \emph{traitors}, in which a legitimate user turns
malicious. Adapting a general-purpose language model to the benign population
and auditing it with sparse features and causal interventions, we find a clean
behavioral detector (top causal features carry essentially no profile mass
despite real psychometrics being present in every input), replicating on real
data the behavioral pole of a profile-versus-behavior dissociation established
on synthetic benchmarks. The detector reaches a mean within-user ROC-AUC of
$0.783$ over $16$ attacked users from mouse, keystroke, and host-event
telemetry alone. But its central finding is a \emph{boundary}: masquerade
windows are detected far better than traitor windows
($0.780$ vs.\ $0.256$ pool ROC-AUC), and---crucially---this gap reproduces
across four unrelated detector families (Isolation Forest, one-class SVM,
PCA reconstruction, and the language model), so it is a property of the
\emph{observation space}, not of any one model: traitors are behaviorally
\emph{calmer} than their own baselines, exfiltrating out of the telemetry
channel. A case study on a near-zero-activity user shows the converse---that
detectability is a property of the serialized field set: enriching the
serialization with a host-event channel flips that user's within-user ROC-AUC
from $0.005$ (inverted) to $0.929$. We argue that reporting what a detector
\emph{structurally cannot} see belongs alongside its detection numbers.
\end{abstract}

\section{Introduction}

An insider-threat detector trained only on benign activity is usually reported
by a single kind of number: an ROC-AUC or a precision at some alert budget,
summarizing how well its anomaly score separates known attacks from normal
behavior. Such numbers answer ``how well does it rank the attacks we gave
it?'' They are silent on two prior questions that matter at least as much in
deployment: \emph{what evidence} is the score actually built on, and \emph{which
kinds of attack are even visible} in the data the detector consumes.

In companion work we developed a mechanistic audit that answers the first
question for adaptation-based sequence detectors, and used it to expose a
dissociation: on one synthetic benchmark a benign-trained language model scores
anomalies through user-identity and profile tokens rather than behavior, while
on another it uses genuine behavioral evidence. That study reserved a real-data
benchmark, TWOS \citep{harilal2017twos}, for a compact confirmation that the
behavioral pole is not a synthetic-data artifact. Here we return to TWOS as the
primary object of study, because it uniquely supports the \emph{second}
question. TWOS contains two operationally different insider scenarios under one
schema---masquerades and traitors---and it pairs high-frequency behavioral
telemetry with each user's real IPIP-50 Big-Five inventory. This lets us ask
not just how well a benign-trained detector scores, but where its blind spots
are and why.

Our contributions are:
\begin{enumerate}[leftmargin=1.3em]
\item \textbf{A real-data audit with real psychometrics.} Adapting Qwen2.5-3B
to the benign TWOS population and running the full sparse-feature-plus-causal
audit, we find the top causally implicated features are essentially $100\%$
behavioral at both audited layers and both adapter seeds---despite a real
personality \texttt{PSY} line in every window. On real data the detector does
not take the identity shortcut, exactly in the population regime
($50\%$ attack prevalence) where identity is a useless separator.
\item \textbf{A model-class-independent masquerade--traitor boundary.} The
language model detects masquerades far better than traitors
($0.780$ vs.\ $0.256$), and this dichotomy holds across Isolation Forest,
one-class SVM, and PCA reconstruction on the same features. Because four
unrelated detector families agree, the boundary is a statement about the
telemetry, not the model: traitor activity is behaviorally quiet.
\item \textbf{A case study localizing detectability to the field set.} A
near-zero-baseline user whose masquerader was the most hyperactive attacker in
the dataset is scored \emph{inverted} (ROC-AUC $0.005$) by a two-source
serialization, and correctly ($0.929$) once a host-event channel is added---a
concrete demonstration that what a detector can see is a property of what the
serialization exposes.
\item \textbf{A causal signature that transfers.} The estimand asymmetry
observed on synthetic data---ablation-dependence is a robust causal signature
of behavioral feature families while patch-repair is configuration-sensitive
---reproduces a fourth time here, on real data.
\end{enumerate}

We take from this that a detector's report should include not only its scores
but an explicit account of what its observation space renders invisible; on
TWOS, that is precisely the traitor.

\section{The TWOS benchmark and serialization}

TWOS \citep{harilal2017twos} records $24$ real users over five working days
across six telemetry sources (keystroke, mouse, host/event-viewer, network,
email, and logon), together with each user's real self-administered IPIP-50
Big-Five personality inventory. Ground truth comprises $12$ staged
\emph{masquerade} sessions ($90$ minutes each, in which an attacker uses a
victim's unlocked machine) and $5$ staged \emph{traitor} sessions ($120$
minutes each, in which a legitimate user acts maliciously). This is the
combination we need: a genuine psychometric axis, two attack modalities, and a
prevalence ($12$ of $24$ users involved in masquerades) that is the opposite
extreme from the small distinctive minority found in the synthetic benchmark
where profile capture occurred.

We serialize the telemetry into the same \texttt{DAY}/\texttt{PSY}/\texttt{SES}
text schema used in the companion audit, at a $10$-minute window grain: a
\texttt{DAY} line with team and machine context, a \texttt{PSY} line carrying
the real scored OCEAN values, and \texttt{SES} lines summarizing per-window
mouse, keystroke, and host-event activity. This yields $18{,}048$ windows,
$138$ positive windows ($101$ masquerade, $37$ traitor), and a $12$-user benign
training pool. We adapt Qwen2.5-3B on the benign windows only (QLoRA, three
epochs, two adapter seeds), score windows by adapted negative log-likelihood,
and run the identical downstream stack (token-delta extraction at layers
$12/18/24$; benign-only top-$k$ sparse autoencoder; full-pool feature
selection; token attribution; causal patching and necessity ablation). No
component of the pipeline is TWOS-specific; the serialization is the only new
code.

\section{The detector is clean and behavioral}

\paragraph{Attribution.} On real data with real personality scores present in
every window, the audit certifies a behavioral detector. The top-five causally
selected features carry essentially all of their attribution mass on
\texttt{SES} (behavioral) tokens and none on \texttt{PSY}/\texttt{DAY}
(profile) tokens: at layer $24$, seed $42$, the five features have profile mass
$[0,0,0,0,0]$ and behavioral mass $\ge 0.987$; the worst single value across
both audited layers and both adapter seeds is $0.983$ behavioral. This is the
fifth clean replication counting the smaller two-source corpus, and it holds
whether or not the enriched host-event channel is included.

This is the outcome the population-structure account predicts. With $12$ of
$24$ users malicious, user identity carries almost no information about which
windows are attacks; a next-token detector therefore has no cheap
identity-shortcut to take, and the behavioral session evidence is what
survives selection. The synthetic benchmark that exhibited profile capture sat
at the opposite corner (four highly distinctive malicious users among
thousands); TWOS confirms the behavioral corner with genuine, rather than
simulated, psychometrics.

\paragraph{Causal support.} Team-matched, same-user-excluded interventions on
the seed-$42$ layer-$24$ top features ($138$ receivers) give a positive
necessity effect ($+0.0019$) and an absent patch-repair effect ($-0.0018$).
The direction reproduces the estimand asymmetry documented three times on CERT:
ablation-dependence is the robust causal signature of behavioral families,
while patch-repair is configuration-sensitive---here plausibly because a
masquerade anomaly is another \emph{human's} behavior, so repairing toward a
benign teammate need not reduce anomaly the way a synthetic profile swap does.
The signature is weak in absolute terms (small pool, thin per-team donor sets)
and we report it as directional corroboration, not as a headline effect.

\section{The masquerade--traitor boundary is a property of the data}

The detector's most consequential result is not a level but a \emph{gap}. Table
\ref{tab:boundary} scores every detector on the same $10$-minute windows, split
by attack modality.

\begin{table}[t]
\centering
\caption{Window ROC-AUC by attack modality on TWOS, across four detector
families trained benign-only on the same serialized features. Every family
detects masquerades far better than traitors; the language model is not an
outlier. ``Within-user'' averages per-user ROC with identity held fixed
(victim's own benign vs.\ attack windows).}
\label{tab:boundary}
\begin{tabular}{lcccc}
\toprule
Detector & All & Masquerade-only & Traitor-only & Mean within-user \\
\midrule
Language model (ours) & $0.639$ & $0.780$ & $0.256$ & $\mathbf{0.783}$ \\
Isolation Forest      & $0.714$ & $0.839$ & $0.374$ & $0.715$ \\
One-class SVM         & $0.607$ & $0.763$ & $0.183$ & $0.626$ \\
PCA reconstruction    & $0.790$ & $0.867$ & $0.582$ & $0.795$ \\
\bottomrule
\end{tabular}
\end{table}

Two things are visible. First, masquerade detection is strong and traitor
detection is at or below chance for \emph{every} detector: the language model
($0.780$ vs.\ $0.256$), Isolation Forest ($0.839$ vs.\ $0.374$), one-class SVM
($0.763$ vs.\ $0.183$), and PCA reconstruction ($0.867$ vs.\ $0.582$) all show
the same ordering with a large margin. Because these four families make
completely different modeling assumptions (density, tree-isolation, margin,
linear reconstruction, and sequence likelihood), their agreement cannot be a
modeling artifact. The boundary is a property of the \emph{observation space}.

Second, the reason is interpretable. A masquerade is a \emph{foreign} operator
on a victim's machine, so the victim's window suddenly looks unlike the
victim---high, out-of-distribution mouse and keystroke activity that every
anomaly model flags. A traitor is the \emph{legitimate} user, and their
malicious phase tends to be behaviorally \emph{calmer} than their own baseline
(disengagement from normal work) while the malicious act itself---data
staged for exfiltration, sent over email---occurs largely \emph{outside} the
mouse/keystroke/host-event channels we serialize. The telemetry that makes
masquerades loud makes traitors quiet. This is not a failure that a better
model fixes; it is a limit of what the field set can express, and any detector
restricted to this observation space inherits it.

The practical consequence is a reporting norm. A benign-trained detector on
this telemetry should not be summarized by its pooled ROC-AUC, which averages a
detectable modality with an undetectable one; it should be reported per
modality, with the traitor result labeled as an observation-space blind spot
rather than a tuning target. The remedy is a channel, not a loss function:
detecting traitors requires serializing the exfiltration surface (email/network
content), not a better anomaly score over mouse and keyboard.

\section{Case study: detectability follows the field set}

If the traitor boundary shows what the observation space \emph{hides}, one user
shows the converse---that enriching the field set changes what a detector can
see. User~21 has a near-empty benign baseline: a median of one mouse event per
$10$-minute window (all quartiles at $1.0$; mean $91$) and roughly $25$
keystrokes---this user barely touched the monitored machine. Their masquerader,
by contrast, was the single most active attacker in TWOS ($23{,}552$ mouse
events per window, against $1$--$2{,}216$ for every other user).

Under a two-source serialization (mouse and keystroke only), this user is
scored \emph{inverted}: within-user ROC-AUC $0.005$, i.e.\ the hyperactive
attack windows are ranked \emph{less} anomalous than the near-empty benign
ones. The mechanism is degeneracy: near-empty benign windows produce
pathological likelihoods that out-rank genuine attack windows, so the raw score
inverts. Adding a host-event channel to the serialization---which gives the
benign windows non-degenerate content and the attack windows a second axis of
evidence---eliminates the inversion entirely: User~21's within-user ROC-AUC
becomes $0.929$.

The lesson generalizes the paper's thesis: the pathological case was not a
model deficiency but a serialization deficiency, and it was fixed by exposing a
field, not by retraining a better scorer. Detectability is a property of what
the serialization renders observable. Reporting a detector therefore means
reporting its field set, because two serializations of the same underlying
telemetry can place the same attack on opposite sides of chance.

\section{Discussion}

Read together, the three results make a single argument about how
benign-trained anomaly detectors should be evaluated and reported. The clean
behavioral audit shows that \emph{what} a detector scores on is not settled by
its score and must be checked. The masquerade--traitor boundary shows that
\emph{which} attacks a detector can see is often not a property of the detector
at all, but of the telemetry it is given---so a good pooled number can hide a
structurally invisible attack class. The User~21 case shows the same coin's
other face: the field set can be changed, and doing so, not tuning the model,
is what moves an attack across the detectability boundary.

None of these are visible at the level of a leaderboard. A detector that scores
$0.78$ pooled on TWOS looks respectable; the audit and the per-modality split
reveal that it is an excellent masquerade detector, a non-detector of traitors,
and a scorer whose behavior on sparse users depends entirely on which channels
were serialized. We suggest that reports of benign-trained insider-threat
detectors carry three artifacts alongside their scores: an attribution of what
evidence the score rests on, a per-modality breakdown that names any
observation-space blind spot, and the field set under which the numbers hold.

\paragraph{Limitations.} TWOS is small ($24$ users, five days, $37$ traitor
windows), so all numbers here are pilot-scale and single-corpus; the causal
interventions in particular rest on thin donor pools and are reported as
directional. The traitor account is a structural hypothesis supported by the
cross-model agreement and by the calmer-than-baseline activity profile, not by
instrumented exfiltration ground truth; confirming it would require serializing
the email and network channels and re-running the split. We regard the boundary
statement---not the individual ROC values---as the transferable claim.

\bibliographystyle{tmlr}
\bibliography{references}

\end{document}
